% !TEX root = ../diploma.tex

\section{Математика}

Теоремы, леммы и определения оформляются окружениями пакета \texttt{amsthm},
их названия задаются в файле \texttt{commands.tex}.

\begin{definition}\label{def:norm}
    \emph{Нормой} вектора $x \in \mathbb{R}^n$ называется величина
    $\|x\| = \sqrt{\sum_{i=1}^{n} x_i^2}$.
\end{definition}

\begin{theorem}[неравенство треугольника]\label{thm:triangle}
    Для любых $x, y \in \mathbb{R}^n$ выполняется $\|x + y\| \le \|x\| + \|y\|$.
\end{theorem}

\begin{proof}
    Следует из неравенства Коши--Буняковского.
\end{proof}

Многострочные формулы удобно набирать окружением \texttt{align}:
\begin{align}
    f(x) &= \int_{0}^{\infty} e^{-\epsilon t} \phi(t) \, dt, \label{eq:laplace} \\
    g(x) &= \sum_{k=1}^{\infty} \frac{\kappa^k}{k!}. \nonumber
\end{align}
Ссылка на формулу: \eqref{eq:laplace}, на теорему: \autoref{thm:triangle}.

Обратите внимание: в \texttt{commands.tex} символы \verb+\epsilon+, \verb+\phi+,
\verb+\kappa+, \verb+\le+ и \verb+\ge+ заменены на варианты, принятые
в русской традиции: $\epsilon$, $\phi$, $\kappa$, $\le$, $\ge$.

Коммутативная диаграмма (пакет \texttt{tikz-cd}):
\[
    \begin{tikzcd}
        A \arrow[r, "f"] \arrow[d, "h"'] & B \arrow[d, "g"] \\
        C \arrow[r, "k"']                & D
    \end{tikzcd}
\]
